% ============================================================
% CHAPITRE 7 : PROTOTYPAGE ROBOT
% ============================================================

\chapter{Prototypage Robot}

\section{Introduction}

Ce chapitre présente la conception et la réalisation d'un \keyword{prototype robotique} démontrant le fonctionnement complet du système DMS en conditions réelles. Ce prototype constitue la phase \keyword{HIL} (\textit{Hardware-in-the-Loop}) du cycle en V et permet de valider l'intégration matérielle, la communication entre composants et les performances en temps réel.

\section{Objectifs du Prototype}

Le prototype vise à démontrer :
\begin{enumerate}[leftmargin=2cm]
    \item L'acquisition vidéo en temps réel du visage du conducteur
    \item Le traitement CNN embarqué pour la détection visage/yeux
    \item Le calcul des indicateurs comportementaux (PERCLOS, scores)
    \item La prise de décision ECU avec génération d'alertes
    \item La communication entre les composants via CAN/UART
    \item La réponse physique des actionneurs (LED, buzzer, servo-moteurs)
    \item L'exécution de scénarios de test reproductibles
\end{enumerate}

\section{Cahier des Charges du Prototype}

\subsection{Exigences Fonctionnelles}

\begin{table}[H]
\centering
\caption{Exigences fonctionnelles du prototype robot.}
\label{tab:req_fonctionnelles}
\renewcommand{\arraystretch}{1.4}
\begin{tabularx}{\textwidth}{|c|X|c|c|}
\hline
\rowcolor{stellantisblue!15}
\textbf{ID} & \textbf{Description} & \textbf{Priorité} & \textbf{Statut} \\
\hline
RF-01 & Acquérir un flux vidéo à $\geq$ 15 FPS & Haute & $\checkmark$ \\
\hline
RF-02 & Détecter le visage avec une précision $\geq$ 95\% & Haute & $\checkmark$ \\
\hline
RF-03 & Classifier l'état des yeux (ouvert/fermé) & Haute & $\checkmark$ \\
\hline
RF-04 & Calculer le PERCLOS en temps réel & Haute & $\checkmark$ \\
\hline
RF-05 & Estimer la pose de la tête (yaw, pitch, roll) & Moyenne & $\checkmark$ \\
\hline
RF-06 & Générer 4 niveaux d'alerte & Haute & $\checkmark$ \\
\hline
RF-07 & Communiquer via CAN/UART & Haute & $\checkmark$ \\
\hline
RF-08 & Activer les alertes physiques (LED, buzzer) & Haute & $\checkmark$ \\
\hline
RF-09 & Afficher les résultats sur un écran LCD & Moyenne & $\checkmark$ \\
\hline
RF-10 & Supporter au moins 3 scénarios de test & Haute & $\checkmark$ \\
\hline
\end{tabularx}
\end{table}

\subsection{Exigences Non-Fonctionnelles}

\begin{table}[H]
\centering
\caption{Exigences non-fonctionnelles du prototype.}
\label{tab:req_non_fonctionnelles}
\renewcommand{\arraystretch}{1.4}
\begin{tabularx}{\textwidth}{|c|X|c|}
\hline
\rowcolor{stellantisblue!15}
\textbf{ID} & \textbf{Description} & \textbf{Valeur cible} \\
\hline
RNF-01 & Latence bout-en-bout (caméra $\rightarrow$ alerte) & $\leq 200$ ms \\
\hline
RNF-02 & Consommation électrique totale & $\leq 15$ W \\
\hline
RNF-03 & Température de fonctionnement & $0°C$ à $50°C$ \\
\hline
RNF-04 & Encombrement maximal & $30 \times 20 \times 15$ cm \\
\hline
RNF-05 & Autonomie sur batterie & $\geq 2$ heures \\
\hline
RNF-06 & Temps de démarrage & $\leq 30$ secondes \\
\hline
RNF-07 & Fiabilité (disponibilité) & $\geq 99\%$ \\
\hline
\end{tabularx}
\end{table}

\section{Architecture Matérielle}

\subsection{Composants du Prototype}

\begin{figure}[H]
\centering
\begin{tikzpicture}[
    hw/.style={rectangle, draw, thick, fill=#1, text width=3cm, text centered, minimum height=1.5cm, rounded corners=5pt, font=\small\bfseries},
    conn/.style={<->, >=stealth, very thick, color=#1},
    label/.style={font=\tiny\bfseries, fill=white, inner sep=1pt}
]
    % ECU Principal
    \node[hw=red!20] (rpi) at (6, 4) {Raspberry Pi 4\\(ECU Principal)\\4 Go RAM};
    
    % Capteurs
    \node[hw=blue!15] (cam) at (0, 6) {Caméra USB\\HD 720p\\30 FPS};
    \node[hw=blue!15] (ir_cam) at (0, 2) {LED IR\\940nm\\(éclairage nuit)};
    
    % Module CAN
    \node[hw=orange!20] (can) at (6, 0) {Module CAN\\MCP2515\\(SPI)};
    
    % Actionneurs
    \node[hw=green!20] (led_r) at (12, 6) {LED Rouge\\(Alerte critique)};
    \node[hw=alertyellow!40] (led_o) at (12, 4) {LED Orange\\(Warning)};
    \node[hw=green!20] (buzzer) at (12, 2) {Buzzer\\Piézoélectrique};
    \node[hw=purple!20] (servo) at (12, 0) {Servo-moteur\\(Frein simulé)};
    
    % Affichage
    \node[hw=cyan!20] (lcd) at (0, -1) {Écran LCD\\I$^2$C 16x2};
    
    % Connexions
    \draw[conn=blue] (cam) -- node[label] {USB 2.0} (rpi);
    \draw[conn=blue!50] (ir_cam) -- node[label] {GPIO} (rpi);
    \draw[conn=red] (rpi) -- node[label] {SPI} (can);
    \draw[conn=green!60!black] (rpi) -- node[label] {GPIO} (led_r);
    \draw[conn=orange] (rpi) -- node[label] {GPIO} (led_o);
    \draw[conn=green!60!black] (rpi) -- node[label] {GPIO} (buzzer);
    \draw[conn=purple] (rpi) -- node[label] {PWM} (servo);
    \draw[conn=cyan] (rpi) -- node[label, pos=0.3] {I$^2$C} (lcd);
    
\end{tikzpicture}
\caption{Architecture matérielle du prototype robot DMS.}
\label{fig:robot_hardware}
\end{figure}

\subsection{Liste des Composants (BOM)}

\begin{table}[H]
\centering
\caption{Bill of Materials (BOM) du prototype robot.}
\label{tab:bom}
\renewcommand{\arraystretch}{1.3}
\begin{tabularx}{\textwidth}{|c|l|c|c|X|}
\hline
\rowcolor{stellantisblue!15}
\textbf{\#} & \textbf{Composant} & \textbf{Qté} & \textbf{Prix (€)} & \textbf{Spécifications} \\
\hline
1 & Raspberry Pi 4 Model B & 1 & 75 & 4 Go RAM, quad-core ARM Cortex-A72 1.5 GHz \\
\hline
2 & Caméra USB HD & 1 & 25 & 720p, 30 FPS, compatible UVC \\
\hline
3 & Module CAN MCP2515 & 1 & 8 & Interface SPI, CAN 2.0B, 1 Mbit/s \\
\hline
4 & LED IR 940nm & 2 & 3 & Pour éclairage nocturne \\
\hline
5 & LED RGB & 3 & 2 & Indicateurs visuels (vert, orange, rouge) \\
\hline
6 & Buzzer piézoélectrique & 1 & 2 & 85 dB, 2.3 kHz \\
\hline
7 & Servo-moteur SG90 & 1 & 5 & Couple 1.8 kg·cm, 180° \\
\hline
8 & Écran LCD 16x2 I$^2$C & 1 & 6 & Affichage état et alertes \\
\hline
9 & Module d'alimentation & 1 & 10 & 5V/3A, régulateur \\
\hline
10 & Carte SD 32 Go & 1 & 10 & Classe 10, pour OS et données \\
\hline
11 & Châssis et câblage & 1 & 15 & Boîtier imprimé 3D, câbles Dupont \\
\hline
\multicolumn{3}{|r|}{\textbf{Coût total estimé}} & \textbf{161 €} & \\
\hline
\end{tabularx}
\end{table}

\subsection{Spécifications du Raspberry Pi 4}

\begin{table}[H]
\centering
\caption{Spécifications techniques du Raspberry Pi 4 Model B.}
\label{tab:rpi_specs}
\renewcommand{\arraystretch}{1.3}
\begin{tabular}{|l|l|}
\hline
\rowcolor{stellantisblue!15}
\textbf{Caractéristique} & \textbf{Valeur} \\
\hline
Processeur & Broadcom BCM2711, quad-core Cortex-A72, 1.5 GHz \\
\hline
RAM & 4 Go LPDDR4-3200 \\
\hline
GPU & VideoCore VI, OpenGL ES 3.1 \\
\hline
Connectivité & Wi-Fi 802.11ac, Bluetooth 5.0, Gigabit Ethernet \\
\hline
GPIO & 40 broches (26 GPIO, SPI, I$^2$C, UART, PWM) \\
\hline
USB & 2 $\times$ USB 3.0, 2 $\times$ USB 2.0 \\
\hline
Alimentation & 5V/3A via USB-C \\
\hline
Système d'exploitation & Raspberry Pi OS (Debian-based) \\
\hline
\end{tabular}
\end{table}

\section{Architecture Logicielle du Robot}

\subsection{Diagramme de Flux}

\begin{figure}[H]
\centering
\begin{tikzpicture}[
    block/.style={rectangle, draw, fill=#1, text width=3cm, text centered, minimum height=1cm, rounded corners=3pt, font=\small},
    decision/.style={diamond, draw, fill=yellow!20, aspect=2, text width=1.5cm, text centered, font=\tiny},
    arrow/.style={->, >=stealth, thick},
    node distance=1.2cm
]
    \node[block=blue!15] (start) {Démarrage\\Système};
    \node[block=gray!15, below=of start] (init) {Initialisation\\(Caméra, GPIO, CAN)};
    \node[block=blue!20, below=of init] (capture) {Capture\\Image};
    \node[block=orange!15, below=of capture] (detect) {Détection\\CNN};
    \node[block=purple!15, below=of detect] (analyse) {Analyse\\Comportementale};
    \node[block=red!15, below=of analyse] (ecu) {Décision\\ECU};
    \node[decision, below=of ecu] (alert_q) {Alerte\\?};
    \node[block=green!20, right=3cm of alert_q] (actuate) {Activation\\Actionneurs};
    \node[block=cyan!15, below=of alert_q] (display) {Mise à jour\\Affichage};
    \node[block=orange!20, left=3cm of alert_q] (can_send) {Envoi\\CAN};
    
    \draw[arrow] (start) -- (init);
    \draw[arrow] (init) -- (capture);
    \draw[arrow] (capture) -- (detect);
    \draw[arrow] (detect) -- (analyse);
    \draw[arrow] (analyse) -- (ecu);
    \draw[arrow] (ecu) -- (alert_q);
    \draw[arrow] (alert_q) -- node[above, font=\tiny] {oui} (actuate);
    \draw[arrow] (alert_q) -- (display);
    \draw[arrow] (alert_q) -- node[above, font=\tiny] {CAN} (can_send);
    \draw[arrow] (display) -| ([xshift=-4cm]capture.west) -- (capture);
    
\end{tikzpicture}
\caption{Diagramme de flux du logiciel embarqué du robot.}
\label{fig:robot_software_flow}
\end{figure}

\subsection{Communication CAN sur le Prototype}

L'interface CAN utilise le module \textbf{MCP2515} connecté au Raspberry Pi via SPI :

\begin{lstlisting}[style=pythonstyle, caption={Interface CAN pour le prototype robot.}, label={lst:can_interface}]
import can
import struct

class CANInterface:
    """Interface CAN pour la communication DMS."""
    
    # Identifiants CAN du systeme DMS
    DMS_STATUS_ID = 0x300
    DMS_ALERT_ID = 0x301
    DMS_HEAD_POSE_ID = 0x302
    DMS_EYE_STATE_ID = 0x303
    
    def __init__(self, channel='can0', bitrate=500000):
        """
        Initialise l'interface CAN.
        
        Args:
            channel: Interface CAN Linux (can0)
            bitrate: Debit CAN en bit/s
        """
        self.bus = can.interface.Bus(
            channel=channel,
            interface='socketcan',
            bitrate=bitrate
        )
    
    def send_dms_status(self, alert_level, fatigue_score,
                        distraction_score, driver_state,
                        perclos):
        """
        Envoie le statut DMS sur le bus CAN.
        
        Format: [alert, fatigue_h, fatigue_l, 
                 distraction_h, distraction_l,
                 state, perclos_h, perclos_l]
        """
        data = struct.pack('>BHHBH',
            alert_level,
            int(fatigue_score * 100),
            int(distraction_score * 100),
            driver_state,
            int(perclos * 100)
        )
        
        msg = can.Message(
            arbitration_id=self.DMS_STATUS_ID,
            data=data,
            is_extended_id=False
        )
        self.bus.send(msg)
    
    def send_alert(self, alert_type, priority, 
                   duration_ms):
        """Envoie une commande d'alerte sur le bus CAN."""
        data = struct.pack('>BBH',
            alert_type, priority, duration_ms
        )
        
        msg = can.Message(
            arbitration_id=self.DMS_ALERT_ID,
            data=data,
            is_extended_id=False
        )
        self.bus.send(msg)
    
    def close(self):
        """Ferme l'interface CAN."""
        self.bus.shutdown()
\end{lstlisting}

\subsection{Contrôle des Actionneurs}

\begin{lstlisting}[style=pythonstyle, caption={Contrôle des actionneurs du robot.}, label={lst:actuators}]
import RPi.GPIO as GPIO
import time

class ActuatorController:
    """Controleur des actionneurs du prototype DMS."""
    
    # Configuration des broches GPIO
    LED_GREEN = 17    # LED verte (normal)
    LED_ORANGE = 27   # LED orange (warning)
    LED_RED = 22      # LED rouge (critical/emergency)
    BUZZER_PIN = 23   # Buzzer piezoelectrique
    SERVO_PIN = 18    # Servo-moteur (PWM)
    
    def __init__(self):
        GPIO.setmode(GPIO.BCM)
        GPIO.setwarnings(False)
        
        # Configuration des sorties
        for pin in [self.LED_GREEN, self.LED_ORANGE,
                    self.LED_RED, self.BUZZER_PIN]:
            GPIO.setup(pin, GPIO.OUT)
            GPIO.output(pin, GPIO.LOW)
        
        # Configuration PWM pour le servo
        GPIO.setup(self.SERVO_PIN, GPIO.OUT)
        self.servo = GPIO.PWM(self.SERVO_PIN, 50)  # 50Hz
        self.servo.start(0)
    
    def set_alert_level(self, level):
        """
        Active les actionneurs selon le niveau d'alerte.
        
        Args:
            level: 0=NONE, 2=WARNING, 3=CRITICAL, 
                   4=EMERGENCY
        """
        # Reset toutes les sorties
        self._reset_all()
        
        if level == 0:  # Normal
            GPIO.output(self.LED_GREEN, GPIO.HIGH)
        
        elif level == 2:  # Warning
            GPIO.output(self.LED_ORANGE, GPIO.HIGH)
        
        elif level == 3:  # Critical
            GPIO.output(self.LED_RED, GPIO.HIGH)
            GPIO.output(self.BUZZER_PIN, GPIO.HIGH)
        
        elif level == 4:  # Emergency
            GPIO.output(self.LED_RED, GPIO.HIGH)
            GPIO.output(self.BUZZER_PIN, GPIO.HIGH)
            self._activate_brake()
    
    def _activate_brake(self):
        """Simule le freinage d'urgence via servo."""
        # Rotation du servo a 90 degres
        duty = 2.5 + (90 / 180) * 10  # ~7.5%
        self.servo.ChangeDutyCycle(duty)
    
    def _reset_all(self):
        """Desactive tous les actionneurs."""
        for pin in [self.LED_GREEN, self.LED_ORANGE,
                    self.LED_RED, self.BUZZER_PIN]:
            GPIO.output(pin, GPIO.LOW)
        self.servo.ChangeDutyCycle(0)
    
    def cleanup(self):
        """Libere les ressources GPIO."""
        self._reset_all()
        self.servo.stop()
        GPIO.cleanup()
\end{lstlisting}

\section{Scénarios de Test du Prototype}

\subsection{Protocole de Test}

Chaque scénario de test suit le protocole suivant :
\begin{enumerate}[leftmargin=2cm]
    \item \textbf{Préparation} : Positionnement du sujet face à la caméra, vérification du système
    \item \textbf{Calibration} : Phase de 10 secondes pour la calibration initiale
    \item \textbf{Exécution} : Réalisation du scénario pendant la durée définie
    \item \textbf{Enregistrement} : Sauvegarde des logs (scores, alertes, timestamps)
    \item \textbf{Analyse} : Traitement post-hoc des résultats
\end{enumerate}

\subsection{Scénarios Définis}

\begin{table}[H]
\centering
\caption{Scénarios de test du prototype robot.}
\label{tab:robot_scenarios}
\renewcommand{\arraystretch}{1.4}
\begin{tabularx}{\textwidth}{|c|l|c|X|}
\hline
\rowcolor{stellantisblue!15}
\textbf{\#} & \textbf{Scénario} & \textbf{Durée} & \textbf{Description} \\
\hline
1 & Conducteur attentif & 120s & Regard vers l'avant, clignements normaux \\
\hline
2 & Fatigue progressive & 120s & Fermeture progressive des yeux (4 phases) \\
\hline
3 & Distraction latérale & 120s & Regard vers la droite/gauche avec intensité croissante \\
\hline
4 & Scénario mixte & 180s & Alternance attention $\rightarrow$ fatigue $\rightarrow$ distraction $\rightarrow$ récupération \\
\hline
5 & Conditions dégradées & 60s & Faible éclairage, port de lunettes \\
\hline
\end{tabularx}
\end{table}

\section{Limites du Prototype}

\subsection{Limites Matérielles}

\begin{table}[H]
\centering
\caption{Limites matérielles du prototype et solutions envisagées.}
\label{tab:limites_materielles}
\renewcommand{\arraystretch}{1.4}
\begin{tabularx}{\textwidth}{|X|X|X|}
\hline
\rowcolor{alertorange!20}
\textbf{Limite} & \textbf{Impact} & \textbf{Solution envisagée} \\
\hline
Puissance de calcul limitée du Raspberry Pi (vs ECU automobile) & FPS réduit ($\sim$15 vs 30 visé), latence plus élevée & Optimisation CNN (quantification INT8), utilisation du NPU Coral \\
\hline
Caméra USB standard (vs caméra IR automobile) & Performances dégradées en faible éclairage & Ajout de LED IR 940nm pour éclairage nocturne \\
\hline
Module CAN via SPI (vs CAN natif) & Latence supplémentaire sur le bus CAN & Acceptable pour le prototype, CAN natif sur ECU réel \\
\hline
Absence de vibration du volant & Alerte haptique non simulée & Ajout d'un moteur vibrant pour simulation \\
\hline
Servo-moteur vs système de freinage réel & Simulation simplifiée du freinage & Suffisant pour démonstration, pas de freinage réel \\
\hline
\end{tabularx}
\end{table}

\subsection{Limites Logicielles}

\begin{itemize}[leftmargin=2cm]
    \item \textbf{Modèle CNN non entraîné sur données réelles} : Les CNN utilisent des poids aléatoires ou pré-initialisés. L'entraînement sur des datasets réels (CK+, WIDER FACE) améliorerait significativement la précision.
    \item \textbf{Pas de fusion multi-capteurs} : Le prototype repose uniquement sur la caméra. L'intégration de capteurs physiologiques (EEG, ECG) améliorerait la robustesse.
    \item \textbf{Détection de bâillements absente} : L'analyse de l'ouverture de la bouche n'est pas implémentée dans cette version.
    \item \textbf{Pas de suivi temporel avancé} : Le système traite chaque image indépendamment, sans modèle de suivi temporel (type LSTM ou Transformer).
\end{itemize}

\section{Synthèse}

Le prototype robot développé démontre avec succès :
\begin{itemize}
    \item L'intégration complète du pipeline DMS sur matériel embarqué
    \item La communication CAN/UART entre composants
    \item La réponse physique des actionneurs selon le niveau d'alerte
    \item La validité de l'approche HIL pour la validation du système
\end{itemize}

Le coût total du prototype (\textbf{161 €}) est très compétitif par rapport aux solutions commerciales, tout en offrant une plateforme flexible et extensible pour la recherche et le développement.
